3.27a与3.26a样品的烧结温度均为1350摄氏度，观察到，3.26a（含有葡萄糖）样品表面炭黑色更深，脆性与碎片化程度更甚，猜测这是由于葡萄糖的碳化留下的游离碳所表现的。相比之下，3.27a样表面呈现可感知的
灰白色，小颗粒下明显有一定硬度，这表明，3.27a样品可能出现了一定程度的碳化硅。然而值得注意的是，在一阶段的实验中，我曾尝试过在1250摄氏度、氩气气氛下烧结出一块具有明显碳化硅特征的样品，
其整体板结，在曲度较大的表面处灰白比凹陷处更甚。两者中，前者有着更高的烧结温度却有着更弱的碳化硅特征，这表明两者的差异并不止在于烧结温度，反逻辑的是，更高温度烧结下，碳化硅特征反而较轻。